\documentclass[../../main.tex]{subfiles}% 注意这里的文件路径不能用 ./main.tex ,否则用latexmk编译子文件会报错
\graphicspath{{\subfix{./image/}}} % 指定图片目录,后续可以直接使用图片文件名
% 注意这里的文件路径不能用 ../../image/ ,否则用latexmk编译子文件会报错

% 例如:
% \begin{figure}[H]
% \centering
% \includegraphics[scale=0.3]{图.png}
% \caption{}
% \label{figure:图}
% \end{figure}
% 注意:上述\label{}一定要放在\caption{}之后,否则引用图片序号会只会显示??.

\begin{document}

\section{向量函数}

\begin{definition}
所谓的\textbf{向量函数}是指从它的定义域到$\mathbb{R}^3$中的映射, 也就是三个有序的实函数. 设有定义在区间$[a,b]$上的向量函数
\begin{align*}
\bm{r}(t) = (x(t), y(t), z(t)),\quad a\leqslant t\leqslant b.
\end{align*}
如果$x(t), y(t), z(t)$都是$t$的连续函数, 则称向量函数$\bm{r}(t)$是\textbf{连续的}; 如果$x(t), y(t), z(t)$都是$t$的连续可微函数, 则称向量函数$\bm{r}(t)$是\textbf{连续可微的}. 向量函数$\bm{r}(t)$的导数和积分的定义与数值函数的导数和积分的定义是相同的, 即
\begin{align*}
\left. \frac{\mathrm{d}\boldsymbol{r}}{\mathrm{d}t} \right|_{t=t_0}&=\lim_{\Delta t\rightarrow 0} \frac{\boldsymbol{r}(t_0+\Delta t)-\boldsymbol{r}(t_0)}{\Delta t}=\lim_{\Delta t\rightarrow 0} \left( \frac{x(t_0+\Delta t)-x(t_0)}{\Delta t},\frac{y(t_0+\Delta t)-y(t_0)}{\Delta t}, \right. \left. \frac{z(t_0+\Delta t)-z(t_0)}{\Delta t} \right) 
\\
&=(x\prime (t_0),y\prime (t_0),z\prime (t_0)),\quad \forall t_0\in (a,b),
\end{align*}
\begin{align*}
\int_a^b{\boldsymbol{r}(t)\mathrm{d}t=\lim_{\lambda \rightarrow 0} \sum_{i=1}^n{\boldsymbol{r}(t_i\prime )\Delta t_i=\left( \int_a^b{x(t)\mathrm{d}t,\int_a^b{y(t)\mathrm{d}t,\int_a^b{z(t)\mathrm{d}t}}} \right) }},
\end{align*}
其中$a = t_0 < t_1 < \cdots < t_n = b$是区间$[a,b]$的任意一个分割, $\Delta t_i = t_i - t_{i-1}$, $t_i'\in[t_{i-1},t_i]$, 并且$\lambda = \max\{\Delta t_i; i=1,\cdots,n\}$. 
\end{definition}
\begin{remark}
这就是说, 向量函数的求导和积分归结为它的分量函数的求导和积分, 因此向量函数的可微性和可积性归结为它的分量函数的可微性和可积性.
\end{remark}

\begin{theorem}
假定$\bm{a}(t),\bm{b}(t),\bm{c}(t)$是三个可微的向量函数, 则它们的内积、向量积和混合积的导数有下面的公式:
\begin{enumerate}[(1)]
\item $(\bm{a}(t)\cdot\bm{b}(t))' = \bm{a}'(t)\cdot\bm{b}(t) + \bm{a}(t)\cdot\bm{b}'(t)$;
\item $(\bm{a}(t)\times\bm{b}(t))' = \bm{a}'(t)\times\bm{b}(t) + \bm{a}(t)\times\bm{b}'(t)$;
\item $(\bm{a}(t),\bm{b}(t),\bm{c}(t))' = (\bm{a}'(t),\bm{b}(t),\bm{c}(t)) + (\bm{a}(t),\bm{b}'(t),\bm{c}(t)) + (\bm{a}(t),\bm{b}(t),\bm{c}'(t))$.
\end{enumerate}
\end{theorem}

\begin{theorem}
设$\bm{a}(t)$是一个处处非零的连续可微的向量函数, 则
\begin{enumerate}[(1)]
\item 向量函数$\bm{a}(t)$的长度是常数当且仅当$\bm{a}'(t)\cdot\bm{a}(t)\equiv 0$.
\item 向量函数$\bm{a}(t)$的方向不变当且仅当$\bm{a}'(t)\times\bm{a}(t)\equiv \bm{0}$.
\item 如果向量函数$\bm{a}(t)$与某一个固定的方向垂直, 那么
\begin{align*}
(\bm{a}(t),\bm{a}'(t),\bm{a}''(t))\equiv 0.
\end{align*}
反过来, 如果上式成立, 并且处处有$\bm{a}'(t)\times\bm{a}(t)\neq 0$, 那么向量函数$\bm{a}(t)$必定与某一个固定的方向垂直.
\end{enumerate}
\end{theorem}

\begin{proof}
\begin{enumerate}[(1)]
\item 因为
\begin{align*}
\frac{\mathrm{d}}{\mathrm{d}t}|\bm{a}(t)|^2 = \frac{\mathrm{d}(\bm{a}(t)\cdot\bm{a}(t))}{\mathrm{d}t} = 2\bm{a}'(t)\cdot\bm{a}(t),
\end{align*}
所以$|\bm{a}(t)|^2$是常数, 当且仅当$\bm{a}'(t)\cdot\bm{a}(t)\equiv 0$.

\item 如果向量函数$\bm{a}(t)$的方向不变, 则有一个固定的单位向量$\bm{b}$, 使得向量函数$\bm{a}(t)$能够写成
\begin{align*}
\bm{a}(t) = f(t)\cdot\bm{b},
\end{align*}
其中$f(t) = \bm{a}(t)\cdot\bm{b}$是处处非零的连续可微函数, 因此
\begin{align*}
\bm{a}'(t) = f'(t)\cdot\bm{b},\quad \bm{a}'(t)\times\bm{a}(t)\equiv 0.
\end{align*}
反过来, 设$\bm{a}'(t)\times\bm{a}(t)\equiv 0$, 命$\bm{b}(t) = \bm{a}(t)/|\bm{a}(t)|$, $|\bm{b}(t)|=1$. 我们要证明$\bm{b}(t)$是常向量函数. 因为$\bm{b}(t)$的长度是$1$, 故由(1)可知,
\begin{align*}
\bm{b}'(t)\cdot\bm{b}(t)\equiv 0,
\end{align*}
即$\bm{b}'(t)\cdot\bm{a}(t)\equiv 0$. 由$\bm{b}(t)$的定义得知
\begin{align*}
\bm{a}(t) = f(t)\bm{b}(t),
\end{align*}
其中$f(t)=|\bm{a}(t)|$处处不为零, 故
\begin{align*}
\bm{a}'(t) &= f'(t)\bm{b}(t) + f(t)\bm{b}'(t), \\
\bm{a}'(t)\times\bm{a}(t) &= f(t)\bm{b}'(t)\times\bm{a}(t) = 0,
\end{align*}
因此$\bm{b}'(t)\times\bm{a}(t)=0$, 即$\bm{b}'(t)$与$\bm{a}(t)$共线, 假设
\begin{align*}
\bm{b}'(t) = \lambda(t)\bm{a}(t).
\end{align*}
由于$\bm{b}'(t)\cdot\bm{a}(t)\equiv 0$, 故
\begin{align*}
\bm{b}'(t)\cdot\bm{a}(t) = \lambda(t)\bm{a}(t)\cdot\bm{a}(t) = \lambda(t)f^2(t)\equiv 0,
\end{align*}
于是$\lambda(t)\equiv 0$, 即
\begin{align*}
\bm{b}'(t)\equiv 0,
\end{align*}
故$\bm{b}(t)$是常向量, 所以向量函数$\bm{a}(t)$的方向不变.

\item 设有单位常向量$\bm{b}$, 使得$\bm{a}(t)\cdot\bm{b}=0$. 对此式求导数得到
\begin{align*}
\bm{a}'(t)\cdot\bm{b}\equiv 0,\quad \bm{a}''(t)\cdot\bm{b}\equiv 0,
\end{align*}
因此向量$\bm{a}(t),\bm{a}'(t),\bm{a}''(t)$都落在$\bm{b}$的正交补空间内, 又因为它们都在三维欧氏空间$E^3$中,所以对于任意的$t$,向量$\bm{a}(t),\bm{a}'(t),\bm{a}''(t)$是共面的, 于是
\begin{align*}
(\bm{a}(t),\bm{a}'(t),\bm{a}''(t))\equiv 0.
\end{align*}
反过来, 假定上面的式子成立, 则
\begin{align*}
(\bm{a}(t)\times\bm{a}'(t))\cdot\bm{a}''(t)\equiv 0.
\end{align*}
已经假设$\bm{a}'(t)\times\bm{a}(t)\neq 0$, 于是可以命
\begin{align*}
\bm{b}(t) = \bm{a}'(t)\times\bm{a}(t),
\end{align*}
则
\begin{align*}
\bm{b}'(t) = \bm{a}''(t)\times\bm{a}(t),
\end{align*}
因而
\begin{align*}
\bm{b}(t)\times\bm{b}'(t) &= \bm{b}(t)\times(\bm{a}''(t)\times\bm{a}(t)) \\
&= (\bm{b}(t)\cdot\bm{a}(t))\bm{a}''(t) - (\bm{b}(t)\cdot\bm{a}''(t))\bm{a}(t) \\
&= (\bm{a}(t),\bm{a}'(t),\bm{a}''(t))\bm{a}(t)\equiv 0,
\end{align*}
根据(2), 向量函数$\bm{b}(t)$有确定的方向. 命$\bm{b}_0 = \frac{\boldsymbol{b}\left( t \right)}{\left| \boldsymbol{b}\left( t \right) \right|}$, 则$\bm{b}_0$是单位常向量, 并且
\begin{align*}
\bm{a}(t)\cdot\bm{b}_0 = \frac{\bm{a}(t)\cdot\bm{b}(t)}{|\bm{b}(t)|}\equiv 0.
\end{align*}
\end{enumerate}

\end{proof}




\end{document}